% Versione italiana del CV ESCO/ATS pubblicato in public/cv.pdf.
% Compilazione: tectonic -o output/pdf cv-source/cv-it.tex
\documentclass[11pt,a4paper]{article}
\usepackage[margin=1.55cm]{geometry}
\usepackage{fontspec}
\setmainfont{lmroman10-regular.otf}[BoldFont=lmroman10-bold.otf,ItalicFont=lmroman10-italic.otf,BoldItalicFont=lmroman10-bolditalic.otf]
\usepackage[italian]{babel}
\usepackage{enumitem}
\usepackage{titlesec}
\usepackage{hyperref}
\hypersetup{hidelinks,pdftitle={Nicolò Ciccarone - Curriculum vitae},pdfauthor={Nicolò Ciccarone},pdflang={it-IT}}
\pagestyle{empty}
\setlength{\parindent}{0pt}
\setlength{\parskip}{0pt}
\setlength{\emergencystretch}{2em}
\titleformat{\section}{\large\bfseries}{}{0pt}{}[\titlerule]
\titlespacing*{\section}{0pt}{12pt}{5pt}
\setlist[itemize]{leftmargin=14pt,itemsep=2pt,topsep=3pt,parsep=0pt}
\newcommand{\job}[4]{\textbf{#1}\hfill #2\\\textit{#3}\hfill #4\par}
\newcommand{\education}[4]{\textbf{#1}\hfill #2\\\textit{#3}\\#4\par\vspace{10pt}}
% Mantieni la gerarchia tramite dimensioni, spaziatura e corsivo, senza grassetto.
\renewcommand{\textbf}[1]{#1}
\let\bfseries\mdseries
\begin{document}
\begin{center}
{\LARGE\bfseries NICOLÒ CICCARONE}\\[3pt]
{\large Senior Data Scientist e statistico applicato | Machine Learning e MLOps}\\[6pt]
{\small Bari, Italia\enspace |\enspace +39 366 322 1036\enspace |\enspace\href{mailto:niciccarone@gmail.com}{niciccarone@gmail.com}\\[2pt]
\href{https://www.linkedin.com/in/nicolociccarone1997/}{linkedin.com/in/nicolociccarone1997}\enspace |\enspace\href{https://github.com/nclcc}{github.com/nclcc}\enspace |\enspace\href{https://nclcc.github.io}{complexalabs.com}}
\end{center}
\vspace{-7pt}
\section*{PROFILO PROFESSIONALE}
Senior data scientist e statistico applicato con formazione in econometria, specializzato nello sviluppo di modelli predittivi e sistemi di machine learning in produzione. Applico inferenza causale, analisi di sopravvivenza e previsione di serie storiche al rischio climatico e finanziario e alle decisioni aziendali. Sviluppo applicazioni di elaborazione dati con Python, SQL, Spark e AWS, integrando API, orchestrazione cloud e monitoraggio. Realizzo soluzioni analitiche per i settori delle risorse umane, del commercio al dettaglio, dell'aviazione e dei servizi finanziari.

\section*{COMPETENZE TECNICHE}
\begin{itemize}
\item \textbf{Modellistica statistica:} statistica, econometria, inferenza causale, modelli di sopravvivenza e mixture-cure, modelli score-driven (GAS), previsione di serie storiche, processi di Hawkes, rischio climatico, gradient boosting.
\item \textbf{Programmazione e dati:} Python, R, SQL, Polars, PySpark / Apache Spark, Julia (base), GIS, dati spaziali, MongoDB; gestione di dati quantitativi.
\item \textbf{Machine Learning:} PyTorch, mixture-of-experts (MoE), elaborazione del linguaggio naturale (NLP), modelli predittivi per il customer lifetime value (CLTV), l'abbandono dei clienti e la previsione della domanda.
\item \textbf{MLOps, cloud e distribuzione:} AWS SageMaker, ECS, EKS, ECR, CloudWatch; Docker, Kubernetes, FastAPI, Celery, Grafana; rilascio di modelli, inferenza in produzione e flussi di lavoro distribuiti.
\item \textbf{Altre piattaforme:} Databricks; Microsoft Azure ML, Pipelines e Fabric; Google Cloud Platform (GCP) Vertex AI Workbench e Colab Enterprise.
\end{itemize}

\section*{ESPERIENZA PROFESSIONALE}
\job{Complexa Labs}{Bari, Italia}{Data Science quantitativa e Data Engineering}{01/2025 - in corso}
\begin{itemize}
\item Modellistica statistica e soluzioni analitiche in produzione per clienti enterprise con fatturato superiore a 1 miliardo di euro.
\item Sviluppo e rilascio su AWS SageMaker di un motore di previsione del tempo di copertura delle posizioni vacanti, basato su modelli di sopravvivenza a tempo discreto score-driven (GAS), pannelli giornalieri con censura a destra e dati del processo di selezione variabili nel tempo.
\item Sviluppo di modelli CLTV con sopravvivenza mixture-cure, regressione del margine tramite gradient boosting e segmentazione calibrata, per prevedere i ricavi e individuare il prodotto successivo da proporre.
\item Progettazione e rilascio di un simulatore di scenari macroeconomici per prevedere gli impatti aziendali al variare delle condizioni di mercato, con orchestrazione distribuita dei flussi di lavoro.
\item Ottimizzazione delle pipeline distribuite con join, funzioni finestra e broadcast in PySpark, \textbf{riducendo del 70\% i tempi di esecuzione}; rilascio di applicazioni NLP come microservizi Python con FastAPI.
\end{itemize}
\vspace{5pt}
\job{acmetric B.V.}{Amsterdam, Paesi Bassi}{Data Scientist}{01/2024 - 11/2024}
\begin{itemize}
\item Sviluppo e rilascio di modelli di intelligenza artificiale causale per l'ottimizzazione dei prezzi e la previsione della domanda presso Adidas International B.V.; l'algoritmo di pricing è tuttora in produzione.
\item Sviluppo e manutenzione di modelli gradient boosting (GBM) per prevedere l'abbandono dei clienti e il contenzioso per LATAM Airlines.
\end{itemize}

\newpage
\job{MOOS Solutions B.V.}{Amsterdam, Paesi Bassi}{Data Scientist}{10/2022 - 12/2023}
\begin{itemize}
\item Progettazione e rilascio di un motore proprietario mixture-of-experts (MoE) per l'ottimizzazione dell'inventario in tempo reale, con instradamento dei dati spaziali di sensori Internet of Things (IoT) verso modelli specializzati; utilizzato da una startup. Valutazione delle previsioni tramite errore percentuale assoluto medio (MAPE).
\end{itemize}
\vspace{6pt}
\job{UBS}{Cracovia, Polonia}{Specialista in modellistica e analisi del rischio}{04/2022 - 09/2022}
\begin{itemize}
\item Sviluppo e rilascio di un primo modello di rischio climatico per un portafoglio immobiliare commerciale, con stress test basati su scenari coerenti con le traiettorie NGFS / IPCC.
\end{itemize}

\section*{ISTRUZIONE E FORMAZIONE SPECIALISTICA}
\education{Università Cattolica del Sacro Cuore, Roma}{04/2026 - 07/2026}{Corso di formazione specialistica in epidemiologia statistica}{Metodi di attribuzione causale delle esposizioni ambientali a esiti di salute pubblica cardiovascolari, respiratori e immunomediati.}
\education{Tinbergen Institute, Amsterdam}{07/2023}{Scuola estiva di livello dottorale in Deep Learning e IA generativa}{Formazione intensiva con Eran Raviv, PhD.}
\education{Vrije Universiteit Amsterdam}{2021 - 2022}{M.Sc. in Econometria e Ricerca Operativa}{Voto: 8,0/10. Tesi: \textit{Price volatility and climatic extremes: distinguishing signal from noise} (Volatilità dei prezzi ed estremi climatici: distinguere il segnale dal rumore).}
\education{Università di Bologna}{2019 - 2021}{Laurea magistrale in Economia, curriculum Econometria}{Voto: 110/110 e lode. Tesi: \textit{Hedging against climate change: A Hawkes process approach} (Copertura contro il cambiamento climatico: un approccio con processi di Hawkes).}
\education{Università di Trento}{2015 - 2018}{Laurea triennale in Economia e Management}{Voto: 110/110 e lode. Honours Programme del Dipartimento di Economia.}
\textbf{Liceo Classico C. Sylos}\\
\textit{Diploma di maturità classica}\\
Voto: 100/100.\par\vspace{10pt}

\section*{LINGUE E RICONOSCIMENTI}
\begin{itemize}
\item \textbf{Lingue:} italiano (madrelingua), inglese (C1/C2), francese (B1), tedesco (A2).
\item \textbf{Riconoscimenti:} borsa di studio per merito INPS (2019); borsa di studio per merito dell'Università di Trento (2018).
\item \textbf{Certificazione:} TOEFL 107/120 (2021).
\end{itemize}
\end{document}
